\section{Conclusion}\label{sec:concl}
A plant needs three things from an identity system: revocation authority shared among its organizations, verification proportional to the consequence of a request, and an enforcement path that does not wait for consensus. We built one around three smart contracts for identity, tiered policy and quorum-governed revocation. Routine requests are decided locally against replicated state, and their decision logs are anchored as Merkle roots; costly requests are decided on-chain with an anomaly gate and a strict mode; batch revocation handles firmware recalls. We proved immediate revocation on the audited path, a $\rho\Delta$ exposure bound on the fast path (conditional on the replica-lag assumption), replay resistance, attestation unforgeability with an expiry time, key-reuse prevention, and governance safety (liveness needs $q$ participating voters). We also stated when a spoofing attacker passes the gate. After the review we also made revocation, reinstatement and recall quorum-only, so a single party can only suspend a device.

The measurements support the design with limits. In simulation, the audited path denied every request after finality, and counted from the submission of the revocation both paths accepted about two requests at one request per second, where token and CRL baselines accepted 33--147 on average. At the two most critical tiers, the gate brought full-spoof acceptance to 0\% and stealthy-mimicry acceptance to 18.8--61.1\% at a cost of 2.2--5.2\% of legitimate requests denied; an ablation shows that the gate (or an equivalent non-linear veto) is what does this, and that strict mode adds nothing once a gate is present. Batch revocation saved about 40\% of gas, and a quorum revocation cost 313k--800k gas for 5--21 voters. With the behavioural signal taken from real traffic, the gate was only as good as the detector behind it. A supervised detector accepted 75--78\% of attack requests of withheld types at the upper tiers of one dataset. A detector trained on normal traffic blocked them but denied 12--22\% of legitimate requests at the two critical tiers, and gates calibrated on legitimate traffic alone let 65--100\% of attacks through for such a detector. An adaptive attacker that sent only its best-scoring flows was accepted 29--77\% of the time on TON\_IoT. An attacker that keeps its behaviour within normal variation is accepted at every tier. A centralised online decision point also matches the zero exposure of the audited path, so the case for this design rests on shared revocation authority and an externally anchored audit trail, not on exposure alone.

On a four-validator Besu QBFT network, audited calls were final in about one 2\,s block. A revocation was visible on a second validator after about 1.9\,s on one host and 1.8\,s across four Google Cloud regions, where it reached a remote validator about 0.2\,s after the submitting node's receipt, and it denied all 30 later audited requests. A three-of-five quorum revocation took 4--8\,s. Sustained throughput was 11--18 decisions per second on one two-vCPU host and at least 31--50 across regions. With two of four validators stopped the chain halted safely, as intended. These are single runs on the contract version before the governance hardening. No round-trip times were recorded, the throughput limit was not determined, and the virtual machines have been deleted, so the runs cannot be repeated. The timing figures of the in-process testbed are compute-bound, and the posture and network signals remain synthetic.

The architecture relies on an external detector and a trust oracle, which it neither provides nor secures. Whether the approach holds up in a real plant has to be settled on a plant network, with real protocols, devices and telemetry. We regard that as the next step, together with threshold-signed oracles, key rotation and voter-set changes (not implemented), and privacy protection for exposed scores. Code, result files and generation scripts are at \href{https://github.com/irfanrazanaqvi/BCZT_IDM_IoT_experments_code}{github.com/irfanrazanaqvi/BCZT\_IDM\_IoT\_experments\_code}; the TON\_IoT and WUSTL data must be obtained from their owners.


\newcommand{\confirm}{\textcolor{red}{[AUTHORS TO CONFIRM]}}
\section*{CRediT authorship contribution statement}
\textbf{Syed Irfan Raza Naqvi:} Conceptualization, Methodology, Software, Validation, Formal analysis, Investigation, Data curation, Writing -- original draft, Visualization \confirm.
\textbf{Muhammad Umar Aftab:} Supervision, Conceptualization, Writing -- review \& editing \confirm.
\textbf{Rana Asif Rehman:} \textcolor{red}{[roles to be completed by the author]}.
\textbf{Ahmad Mohsin:} \textcolor{red}{[roles to be completed by the author]}.
\textbf{Sayed Qaiser Ali Shah:} \textcolor{red}{[roles to be completed by the author]}.

\section*{Declaration of competing interest}
The authors declare that they have no known competing financial interests or personal relationships that could have appeared to influence the work reported in this paper. \confirm

\section*{Funding}
This research did not receive any specific grant from funding agencies in the public, commercial, or not-for-profit sectors. \confirm

\section*{Data availability}
The contracts, simulation and experiment scripts, and all result files are openly available at \href{https://github.com/irfanrazanaqvi/BCZT_IDM_IoT_experments_code}{github.com/irfanrazanaqvi/BCZT\_IDM\_IoT\_experments\_code}. The TON\_IoT and WUSTL-IIoT datasets are third-party resources and must be obtained from their owners under their licences.

\section*{Declaration of generative AI and AI-assisted technologies in the writing process}
During the preparation of this work the authors used Claude (Anthropic) to assist with drafting and editing text, generating and debugging experiment and figure code, and organising the reference list. After using this tool, the authors reviewed and edited the content as needed and take full responsibility for the content of the publication. \confirm


\begingroup\small\setstretch{1.0}
\begin{thebibliography}{99}\small
\setlength{\itemsep}{0pt}
\setlength{\parskip}{0pt}
\small
\setlength{\itemsep}{0pt}
\setlength{\parskip}{0pt}
\bibitem{stouffer2023sp80082}
K. Stouffer \emph{et al.}, ``Guide to Operational Technology (OT) Security,'' NIST Special Publication 800-82 Rev.~3, National Institute of Standards and Technology, Sep. 2023. \href{https://doi.org/10.6028/NIST.SP.800-82r3}{doi:\nolinkurl{10.6028/NIST.SP.800-82r3}}

\bibitem{nist800207}
S. Rose, O. Borchert, S. Mitchell, and S. Connelly, ``Zero Trust Architecture,'' NIST Special Publication 800-207, National Institute of Standards and Technology, 2020. \href{https://doi.org/10.6028/NIST.SP.800-207}{doi:\nolinkurl{10.6028/NIST.SP.800-207}}

\bibitem{fernandez2024critical}
E. B. Fernandez and A. Brazhuk, ``A critical analysis of zero trust architecture (ZTA),'' \emph{Computer Standards \& Interfaces}, vol. 89, Art. no. 103832, 2024. \href{https://doi.org/10.1016/j.csi.2024.103832}{doi:\nolinkurl{10.1016/j.csi.2024.103832}}

\bibitem{phiayura2023comprehensive}
P. Phiayura and S. Teerakanok, ``A comprehensive framework for migrating to zero trust architecture,'' \emph{IEEE Access}, vol. 11, pp. 19487--19511, 2023. \href{https://doi.org/10.1109/ACCESS.2023.3248622}{doi:\nolinkurl{10.1109/ACCESS.2023.3248622}}

\bibitem{song2025ztot}
M. W. Song, T. Nguyen, and C. Irvine, ``A zero trust architecture for critical operational technology systems,'' in \emph{Proc. 58th Hawaii Int. Conf. Syst. Sci. (HICSS)}, 2025. \href{https://doi.org/10.24251/HICSS.2025.849}{doi:\nolinkurl{10.24251/HICSS.2025.849}}

\bibitem{zanasi2024flexible}
C. Zanasi, S. Russo, and M. Colajanni, ``Flexible zero trust architecture for the cybersecurity of industrial IoT infrastructures,'' \emph{Ad Hoc Networks}, vol. 156, Art. no. 103414, 2024. \href{https://doi.org/10.1016/j.adhoc.2024.103414}{doi:\nolinkurl{10.1016/j.adhoc.2024.103414}}

\bibitem{w3cdid2022}
M. Sporny \emph{et al.}, ``Decentralized Identifiers (DIDs) v1.0,'' W3C Recommendation, Jul. 2022. \href{https://www.w3.org/TR/did-core/}{w3.org/TR/did-core}

\bibitem{hosseini2023blockchain}
S. M. Hosseini, J. Ferreira, and P. C. Bartolomeu, ``Blockchain-based decentralized identification in IoT: an overview of existing frameworks and their limitations,'' \emph{Electronics}, vol. 12, no. 6, p. 1283, 2023. \href{https://doi.org/10.3390/electronics12061283}{doi:\nolinkurl{10.3390/electronics12061283}}

\bibitem{aguru2025ssi}
A. D. Aguru and S. B. Erukala, ``A blockchain-based self-sovereign identity management framework for Internet of Things,'' in \emph{Recent Trends in Intelligent Systems and Next Generation Wireless Communication (IIWCS 2024)}, Lecture Notes in Networks and Systems, vol. 1329. Singapore: Springer, 2025, pp. 299--309. \href{https://doi.org/10.1007/978-981-96-4741-5_26}{doi:\nolinkurl{10.1007/978-981-96-4741-5_26}}

\bibitem{folino2025puf}
G. Folino, A. Forestiero, and G. Papuzzo, ``Self-sovereign identification of IoT devices by using physically unclonable functions and blockchain,'' in \emph{Numerical Computations: Theory and Algorithms (NUMTA 2023)}, Lecture Notes in Computer Science, vol. 14478. Cham: Springer, 2025. \href{https://doi.org/10.1007/978-3-031-81247-7_26}{doi:\nolinkurl{10.1007/978-3-031-81247-7_26}}

\bibitem{zhai2025chebyshev}
Z. Zhai \emph{et al.}, ``A lightweight authentication method for Industrial Internet of Things based on blockchain and Chebyshev chaotic maps,'' \emph{Future Internet}, vol. 17, no. 8, Art. no. 338, 2025. \href{https://doi.org/10.3390/fi17080338}{doi:\nolinkurl{10.3390/fi17080338}}

\bibitem{essaid2025iiotbc}
M. Essaid and H. Ju, ``Blockchain solutions for enhancing security and privacy in Industrial IoT,'' \emph{Applied Sciences}, vol. 15, no. 12, Art. no. 6835, 2025. \href{https://doi.org/10.3390/app15126835}{doi:\nolinkurl{10.3390/app15126835}}

\bibitem{liu2024dissecting}
C. Liu, R. Tan, Y. Wu, Y. Feng, Z. Jin, F. Zhang, Y. Liu, and Q. Liu, ``Dissecting zero trust: research landscape and its implementation in IoT,'' \emph{Cybersecurity}, vol. 7, no. 1, Art. no. 20, 2024. \href{https://doi.org/10.1186/s42400-024-00212-0}{doi:\nolinkurl{10.1186/s42400-024-00212-0}}

\bibitem{nie2025zeroTrust}
S. Nie, J. Ren, R. Wu, P. Han, Z. Han, and W. Wan, ``Zero-trust access control mechanism based on blockchain and inner-product encryption in the Internet of Things in a 6G environment,'' \emph{Sensors}, vol. 25, no. 2, p. 550, 2025. \href{https://doi.org/10.3390/s25020550}{doi:\nolinkurl{10.3390/s25020550}}

\bibitem{salim2026ztblociot}
M. M. Salim, M. Kim, S. K. Singh, and J. H. Park, ``Zero-trust blockchain-enabled framework for scalable and secure IoT networks (ZT-BlocIoT),'' \emph{Future Generation Computer Systems}, vol. 175, Art. no. 108093, 2026. \href{https://doi.org/10.1016/j.future.2025.108093}{doi:\nolinkurl{10.1016/j.future.2025.108093}}

\bibitem{mukta2026ztdelegation}
R. Mukta, S. Pal, K. Chowdhury, M. Hitchens, H. Paik, and S. S. Kanhere, ``Zero trust-driven access control delegation using blockchain,'' \emph{Blockchain: Research and Applications}, vol. 7, no. 1, Art. no. 100319, 2026. \href{https://doi.org/10.1016/j.bcra.2025.100319}{doi:\nolinkurl{10.1016/j.bcra.2025.100319}}

\bibitem{jamwal2026survey}
E. Jamwal, P. Srivastava, and P. S. Chakraborty, ``Blockchain-enabled zero-trust framework for secure cloud data sharing: survey,'' in \emph{Proc. 2026 Int. Conf. Intell. Syst. Eng., Secured Syst. Cybersecurity (ICISESSC)}, 2026. \href{https://doi.org/10.1109/ICISESSC68634.2026.11542695}{doi:\nolinkurl{10.1109/ICISESSC68634.2026.11542695}}

\bibitem{farhan2026ztfsf}
M. Farhan \emph{et al.}, ``ZTF-SF: a zero-trust blockchain-enabled IIoT framework for secure smart factories,'' in \emph{Proc. 2026 2nd Int. Conf. Comput. Intell. Approaches Appl. (ICCIAA)}, 2026. \href{https://doi.org/10.1109/ICCIAA68481.2026.11543806}{doi:\nolinkurl{10.1109/ICCIAA68481.2026.11543806}}

\bibitem{potturi2026hybrid}
K. Potturi, B. Gautam, C. S. Pareek, A. Gupta, S. C. Chava, and B. S. Katta, ``A hybrid AI and blockchain-enabled zero-trust security framework for cloud, edge, and IoT systems,'' in \emph{Proc. 2026 Int. Conf. Comput. Theory Wireless Commun. (ICCTWC)}, 2026. \href{https://doi.org/10.1109/ICCTWC68241.2026.11582964}{doi:\nolinkurl{10.1109/ICCTWC68241.2026.11582964}}

\bibitem{wang2026medicaliot}
Z. Wang, Y. Guo, and Y. Guo, ``Blockchain-based zero-trust continuous authentication scheme for medical IoT,'' \emph{Journal of Parallel and Distributed Computing}, vol. 217, Art. no. 105333, 2026. \href{https://doi.org/10.1016/j.jpdc.2026.105333}{doi:\nolinkurl{10.1016/j.jpdc.2026.105333}}

\bibitem{james2024authz}
M. James, T. Newe, D. O'Shea, and G. D. O'Mahony, ``Authentication and authorization in zero trust IoT: A survey,'' in \emph{Proc. Irish Signals and Systems Conf. (ISSC)}, 2024. \href{https://doi.org/10.1109/ISSC61953.2024.10603175}{doi:\nolinkurl{10.1109/ISSC61953.2024.10603175}}

\bibitem{kumar2024industry40}
R. Kumar and R. Sharma, ``A comprehensive approach to industry 4.0 security: blockchain and dynamic access control integration,'' in \emph{Proc. 2024 6th Int. Conf. Computational Intelligence and Communication Technologies (CCICT)}, 2024, pp. 477--484. \href{https://doi.org/10.1109/CCICT62777.2024.00080}{doi:\nolinkurl{10.1109/CCICT62777.2024.00080}}

\bibitem{azad2024verify}
M. A. Azad, S. Abdullah, J. Arshad, H. Lallie, and Y. H. Ahmed, ``Verify and trust: a multidimensional survey of zero-trust security in the age of IoT,'' \emph{Internet of Things}, vol. 27, Art. no. 101227, 2024. \href{https://doi.org/10.1016/j.iot.2024.101227}{doi:\nolinkurl{10.1016/j.iot.2024.101227}}

\bibitem{nahar2024survey}
N. Nahar, K. Andersson, O. Schel\'en, and S. Saguna, ``A survey on zero trust architecture: applications and challenges of 6G networks,'' \emph{IEEE Access}, vol. 12, pp. 94753--94764, 2024. \href{https://doi.org/10.1109/ACCESS.2024.3425350}{doi:\nolinkurl{10.1109/ACCESS.2024.3425350}}

\bibitem{prezta2026}
Z. Wei, O. M. Ameer, Y. Zhang, and N. Borisov, ``PREZTA: Provable remote execution of zero-trust authorization using SNARKs,'' arXiv:2607.11466, 2026. \href{https://doi.org/10.48550/arXiv.2607.11466}{doi:\nolinkurl{10.48550/arXiv.2607.11466}}

\bibitem{ren2025aiot}
J. Ren, J. Zhang, Y. Ren, and J. Xu, ``Blockchain-based self-sovereign identity management mechanism in AIoT environments,'' \emph{Electronics}, vol. 14, no. 19, p. 3954, 2025. \href{https://doi.org/10.3390/electronics14193954}{doi:\nolinkurl{10.3390/electronics14193954}}

\bibitem{ali2025dtaccess}
G. Ali, S. Shah, M. Elaffendi, and N. Ahmad, ``Blockchain-based access management framework for interoperable digital twins in industrial IoT,'' \emph{Frontiers in Blockchain}, vol. 8, Art. no. 1693926, 2025. \href{https://doi.org/10.3389/fbloc.2025.1693926}{doi:\nolinkurl{10.3389/fbloc.2025.1693926}}

\bibitem{philipp2023daxiot}
A. Philipp and A. K\"upper, ``DAXiot: A decentralized authentication and authorization scheme for dynamic IoT networks,'' arXiv:2307.06919, 2023. \href{https://doi.org/10.48550/arXiv.2307.06919}{doi:\nolinkurl{10.48550/arXiv.2307.06919}}

\bibitem{mazzocca2024evoke}
C. Mazzocca, A. Acar, S. Uluagac, and R. Montanari, ``EVOKE: Efficient revocation of verifiable credentials in IoT networks,'' in \emph{Proc. 33rd USENIX Security Symp.}, Philadelphia, PA, USA, Aug. 2024, pp. 1279--1295. \href{https://www.usenix.org/system/files/usenixsecurity24-mazzocca.pdf}{usenix.org (PDF)}

\bibitem{kumar2026privacy}
A. Kumar, B. Sharma, R. K. Kaushal, and I. Ben Dhaou, ``Privacy-preserving identity management system leveraging blockchain technology,'' in \emph{Proc. 15th IEEE Int. Conf. Commun. Syst. Netw. Technol. (CSNT)}, 2026, pp. 625--631. \href{https://doi.org/10.1109/CSNT69054.2026.11502058}{doi:\nolinkurl{10.1109/CSNT69054.2026.11502058}}

\bibitem{rivera2026yangbaxter}
Y. E. Rivera-Julio, E. A. Leal-Narv\'aez, and J. Prieto Tejedor, ``Blockchain-backed revocation and Yang--Baxter consistency screening for zero-trust IoT admission control,'' \emph{Sensors}, vol. 26, no. 18, Art. no. 5960, 2026. \href{https://doi.org/10.3390/s26185960}{doi:\nolinkurl{10.3390/s26185960}}

\bibitem{usama2026multilayer}
M. Usama, A. Aziz, N. Alasbali, N. Alturki, M. Hanif, and M. Ur Rehman, ``Blockchain-enabled identity management for IoT: a multi-layered defense against adversarial AI,'' \emph{Scientific Reports}, vol. 16, Art. no. 4371, 2026. \href{https://doi.org/10.1038/s41598-026-35208-y}{doi:\nolinkurl{10.1038/s41598-026-35208-y}}

\bibitem{tom2025flcti}
A. K. Tom, A. Khraisat, T. Jan, M. Whaiduzzaman, T. D. Nguyen, and A. Alazab, ``Survey of federated learning for cyber threat intelligence in Industrial IoT: Techniques, applications and deployment models,'' \emph{Future Internet}, vol. 17, no. 9, Art. no. 409, 2025. \href{https://doi.org/10.3390/fi17090409}{doi:\nolinkurl{10.3390/fi17090409}}

\bibitem{shujaa2025enhancing}
W. Shujaa, M. Alanzi, and S. Sankaranarayanan, ``Enhancing IoT security through blockchain integration,'' \emph{Frontiers in Computer Science}, vol. 7, Art. no. 1670473, 2025. \href{https://doi.org/10.3389/fcomp.2025.1670473}{doi:\nolinkurl{10.3389/fcomp.2025.1670473}}

\bibitem{hassan2023oracle}
A. Hassan, I. Makhdoom, W. Iqbal, A. Ahmad, and A. Raza, ``From trust to truth: Advancements in mitigating the blockchain oracle problem,'' \emph{J. Netw. Comput. Appl.}, vol. 217, Art. no. 103672, 2023. \href{https://doi.org/10.1016/j.jnca.2023.103672}{doi:\nolinkurl{10.1016/j.jnca.2023.103672}}

\bibitem{wei2024quality}
Z. Wei \emph{et al.}, ``Survey on quality assurance of smart contracts,'' \emph{ACM Comput. Surv.}, vol. 57, no. 2, Art. no. 32, 2024. \href{https://doi.org/10.1145/3695864}{doi:\nolinkurl{10.1145/3695864}}

\bibitem{ta2024gascost}
M. T. Ta and T. Q. Do, ``A study on gas cost of Ethereum smart contracts and performance of blockchain on simulation tool,'' \emph{Peer-to-Peer Networking and Applications}, vol. 17, no. 1, pp. 200--212, 2024. \href{https://doi.org/10.1007/s12083-023-01598-3}{doi:\nolinkurl{10.1007/s12083-023-01598-3}}

\bibitem{hasan2023smartcontract}
M. R. Hasan, A. Alazab, S. B. Joy, M. N. Uddin, M. A. Uddin, A. Khraisat, I. Gondal, W. F. Urmi, and M. A. Talukder, ``Smart contract-based access control framework for Internet of Things devices,'' \emph{Computers}, vol. 12, no. 11, Art. no. 240, 2023. \href{https://doi.org/10.3390/computers12110240}{doi:\nolinkurl{10.3390/computers12110240}}

\bibitem{haque2024scalable}
E. U. Haque, A. Shah, J. Iqbal, S. S. Ullah, R. Alroobaea, and S. Hussain, ``A scalable blockchain based framework for efficient IoT data management using lightweight consensus,'' \emph{Scientific Reports}, vol. 14, Art. no. 7841, 2024. \href{https://doi.org/10.1038/s41598-024-58578-7}{doi:\nolinkurl{10.1038/s41598-024-58578-7}}

\bibitem{peepliwal2024ztab}
A. K. Peepliwal, H. M. Pandey, S. Prakash, A. A. Mahajan, S. S. Chowhan, V. Kumar, and R. Sharma, ``A prototype model of zero trust architecture blockchain with EigenTrust-based practical Byzantine fault tolerance protocol to manage decentralized clinical trials,'' \emph{Digital Communications and Networks}, 2024. \href{https://www.sciencedirect.com/science/article/pii/S2096720924000459}{sciencedirect.com (S2096720924000459)}

\bibitem{madhwal2025gasless}
Y. Madhwal \emph{et al.}, ``Empowering autonomous IoT devices in blockchain through gasless transactions,'' \emph{Blockchain: Research and Applications}, vol. 6, no. 2, Art. no. 100257, 2025. \href{https://doi.org/10.1016/j.bcra.2024.100257}{doi:\nolinkurl{10.1016/j.bcra.2024.100257}}

\bibitem{ren2026ditdp}
L. Ren, H. Li, Y. Liu, D. Liao, M. Zhang, and H. Jin, ``Identity trust management based on differential privacy for internet of vehicles,'' \emph{Scientific Reports}, vol. 16, Art. no. 20921, 2026. \href{https://doi.org/10.1038/s41598-026-50273-z}{doi:\nolinkurl{10.1038/s41598-026-50273-z}}

\bibitem{alsaedi2020ton}
A. Alsaedi, N. Moustafa, Z. Tari, A. Mahmood, and A. Anwar, ``TON\_IoT telemetry dataset: A new generation dataset of IoT and IIoT for data-driven intrusion detection systems,'' \emph{IEEE Access}, vol. 8, pp. 165130--165150, 2020. \href{https://doi.org/10.1109/ACCESS.2020.3022862}{doi:\nolinkurl{10.1109/ACCESS.2020.3022862}}

\bibitem{moustafa2021network}
N. Moustafa, ``A new distributed architecture for evaluating AI-based security systems at the edge: Network TON\_IoT datasets,'' \emph{Sustainable Cities and Society}, vol. 72, Art. no. 102994, 2021. \href{https://doi.org/10.1016/j.scs.2021.102994}{doi:\nolinkurl{10.1016/j.scs.2021.102994}}

\bibitem{booij2022ton}
T. M. Booij, I. Chiscop, E. Meeuwissen, N. Moustafa, and F. T. H. den Hartog, ``ToN\_IoT: The role of heterogeneity and the need for standardization of features and attack types in IoT network intrusion data sets,'' \emph{IEEE Internet Things J.}, vol. 9, no. 1, pp. 485--496, 2022. \href{https://doi.org/10.1109/JIOT.2021.3085194}{doi:\nolinkurl{10.1109/JIOT.2021.3085194}}

\bibitem{teixeira2018scada}
M. A. Teixeira, T. Salman, M. Zolanvari, R. Jain, N. Meskin, and M. Samaka, ``SCADA system testbed for cybersecurity research using machine learning approach,'' \emph{Future Internet}, vol. 10, no. 8, Art. no. 76, 2018. \href{https://doi.org/10.3390/fi10080076}{doi:\nolinkurl{10.3390/fi10080076}}

\bibitem{teixeira2018wustl}
M. A. Teixeira, M. Zolanvari, and R. Jain, ``WUSTL-IIOT-2018 dataset for ICS (SCADA) cybersecurity research,'' IEEE DataPort, 2018. \href{https://doi.org/10.21227/kzgp-7t84}{doi:\nolinkurl{10.21227/kzgp-7t84}}

\bibitem{zolanvari2019vulnerability}
M. Zolanvari, M. A. Teixeira, L. Gupta, K. M. Khan, and R. Jain, ``Machine learning-based network vulnerability analysis of industrial Internet of Things,'' \emph{IEEE Internet Things J.}, vol. 6, no. 4, pp. 6822--6834, 2019. \href{https://doi.org/10.1109/JIOT.2019.2912022}{doi:\nolinkurl{10.1109/JIOT.2019.2912022}}

\bibitem{zolanvari2021wustl}
M. Zolanvari, ``WUSTL-IIOT-2021,'' IEEE DataPort, 2021. \href{https://doi.org/10.21227/yftq-n229}{doi:\nolinkurl{10.21227/yftq-n229}}

\end{thebibliography}
\endgroup
\end{document}
